% RESEARCH_DRAFT_SHA256: 078479004706d2dca88f4f0daaa6a013c58720b3f873cba2ef77f7e800e04942
% Layout and source-reference wrappers expanded for the web renderer.
\documentclass[11pt]{article}
\usepackage[margin=25mm]{geometry}
\usepackage{fontspec,unicode-math}
\setmainfont{Latin Modern Roman}
\setmathfont{Latin Modern Math}
\usepackage{amsmath,amssymb,mathtools,xurl,hyperref}
\hypersetup{colorlinks=true,urlcolor=blue,linkcolor=blue}
\setlength{\parindent}{0pt}
\setlength{\parskip}{5pt}
\setlength{\emergencystretch}{3em}
\newcommand{\sref}[2]{\hyperlink{source:#1:#2}{[S#2]}}
\newcommand{\cataloguescope}[1]{\paragraph{Source question.}#1}
\newcommand{\floor}[1]{\left\lfloor#1\right\rfloor}
\begin{document}

\section*{Two-complex embeddings: a local obstruction and the certificate-size gap}
\noindent UnsolvedMath ID 30002956; OWR-13940-003. Research attempt dated 11 October 2026.
\paragraph{Result in scope.}This manuscript contains a source-specific attempt and complete proofs of the stated auxiliary calculations. It makes no new solution or novelty claim for the entire target.
\section{Original problem and preserved definition}
The following definition and source list are reproduced from the repository. Any earlier inspection dates and reported literature claims in this preserved text describe that source record, not additional verification by this attempt.
% BEGIN PRESERVED DEFINITION
\subsection{Definitions and mathematical statement}

\paragraph{Shared notation.}Standard number systems are denoted by $\mathbb N,\mathbb Z,\mathbb Q,\mathbb R,\mathbb C$; index ranges, computational models and specialized notation follow the cited problem source and the definitions below.


A finite simplicial complex is a finite family of vertex subsets closed under taking subsets. Its dimension is one less than its largest simplex size. A piecewise-linear embedding is an injective map linear on simplices after finite subdivisions.
\paragraph{Statement recorded in the supplied source.}
Determine the decidability and computational complexity of deciding whether a given 2-complex embeds in a 3-manifold, including the case of non-orientable surfaces of even Euler genus and the related thickenability problem. \href{https://doi.org/10.4171/owr/2015/45}{[S2]}
\paragraph{Scope qualification.}
The ambient manifold representation, fixed versus input manifold, and relationship of the surface and thickenability subquestions are not fully specified in the short question.
\paragraph{Source status.}
The supplied catalogue labels this record partially solved. Its source-status metadata was checked against the supplied snapshot on 2026-10-11. This import does not independently establish current open status or mathematical correctness. The catalogue literature assessment is dated 2026-08-21.
\subsection{Short English statement}
Classify decidability and complexity of embedding the specified two-dimensional complexes in three-manifolds.
\subsection{Sources}
\begin{itemize}
\item[S1] Ulam AI and the UnsolvedMath Contributors, UnsolvedMath: A Curated Collection of Open Mathematics Problems, record 30002956 (OWR-13940-003). Catalogue curation and metadata: CC BY 4.0. The supplied source question is retained; definitions and the English summary are editorial additions.
\url{https://huggingface.co/datasets/ulamai/UnsolvedMath}.
\item[S2] Computational Geometric and Algebraic Topology (2016), Oberwolfach Reports, DOI 10.4171/owr/2015/45, pp. 2637--2699
\url{https://doi.org/10.4171/owr/2015/45}.
\item[S3] Burton--de Mesmay--Wagner, Finding Non-orientable Surfaces
\url{https://doi.org/10.1007/s00454-017-9900-0}.
\item[S4] Matoušek--Sedgwick--Tancer--Wagner, Embeddability in the 3-sphere
\url{https://arxiv.org/abs/1402.0815}.
\textit{Underlying Oberwolfach Reports material retains its source-specific licence, including CC BY-SA 4.0 where applicable; see THIRD\_PARTY\_NOTICES.md.}
\end{itemize}
% END PRESERVED DEFINITION
\subsection{Verbatim supplied catalogue question}
Determine the decidability and computational complexity of deciding whether a given 2-complex embeds in a 3-manifold, including the case of non-orientable surfaces of even Euler genus and the related thickenability problem.

\section{The original distinctions}
De Mesmay's report, printed pages 2666--2668, separates embeddings into a
specified 3-manifold, embeddings into $\mathbb R^3$, existence of some
3-manifold containing the complex (thickenability), and non-orientable surface
subcases. The odd-genus result in the report does not cover even genus.
We keep the piecewise-linear convention in the preserved definition and test
an attempted route based on contractibility, local links, and finite
simplicial embedding certificates. No missing ambient-manifold quantifier is
filled in by silently choosing the easiest of these problems.

\section{A contractible complex that fails every ambient choice}
Let $C$ be the simplicial cone over the graph $K_5$, with apex $v$.
It is contractible: in cone coordinates $(x,t)$, where all $(x,0)$ are the
apex, the homotopy $(x,t)\mapsto(x,(1-s)t)$ contracts it as $s$ goes from
zero to one. This is well defined at the collapsed apex and continuous on
each cone simplex, hence on the finite complex.
Nevertheless $C$ has no PL embedding into any 3-manifold.

We justify the local obstruction. In a PL coordinate neighborhood of the
image of $v$, subdivide so the embedding is linear on the finitely many
incident simplices. A sufficiently small sphere centered at that image
meets each incident edge in a point and each incident triangular sector in
an arc. The arcs are disjoint except at their prescribed endpoints because
the map is injective. Equivalently, radial intersection embeds a subdivision
of the link of $v$ in the surrounding two-sphere. Shrinking the sphere keeps
all nonincident simplices away: there are finitely many compact nonincident
pieces disjoint from the apex. If the image lies on a manifold boundary,
the local sphere is a hemisphere, which still embeds in a sphere. Thus the
vertex link must be a planar graph. At the cone apex it is exactly $K_5$;
subdivision cannot cure a nonplanar graph, since suppressing a degree-two
vertex simply concatenates its two embedded arcs.

For completeness $K_5$ is nonplanar by the elementary planar edge bound.
A connected simple planar graph with at least three vertices and no bridges
has Euler relation $v-e+f=2$ and every face boundary has length at least
three. Counting edge sides gives $3f\le2e$, hence $e\le3v-6$.
Euler's relation follows by deleting cycle edges one at a time: deleting a
nonbridge merges its two incident faces, decreasing both $e$ and $f$ by one;
what remains is a spanning tree, with $e=v-1$ and one face.
For $K_5$, $v=5,e=10>9$. This proves the contradiction.
The cone example is explicitly mentioned by the primary report; the proof
here rederives it rather than claiming a new obstruction.

This also defeats the proposed inference that contraction to a point should
permit an embedding into a 3-ball. A contraction is a continuous family of
maps; its intermediate maps need not be injective. It provides no control
of the nonplanar local link. Checking global homotopy triviality therefore
cannot replace local embedding constraints.

\section{What an explicit simplicial certificate does certify}
For fixed finite complexes $K'$ and $M'$, a map on their vertices that is
injective and sends each simplex of $K'$ to a simplex of $M'$ gives an
embedding of their realizations. To prove this, the affine map is injective
on each simplex because its image vertices form an affinely independent
simplex. Images of two source simplices intersect exactly in the face on
their common image vertices, by the defining intersection property of an
ambient simplicial complex. Vertex injectivity makes these precisely the
images of the common source vertices. Thus the map is globally injective.
It is continuous, and its inverse onto the image is continuous because a
continuous bijection from a compact space to a Hausdorff space is closed:
a closed subset is compact and its image is compact, hence closed.
This proves a topological, and indeed simplicial, embedding.

Such a certificate is checkable in polynomial time in the explicit sizes
of $K',M'$: compare vertex images for injectivity and check each simplex
image against the target simplex list. But the original inputs are $K,M$,
not arbitrary already-supplied refinements. An attempted exhaustive search
would need certificates that $K',M'$ are the appropriate subdivisions and
would need a bound on their sizes to prove NP membership or a running-time
classification. Finite enumerations through a bounded set of subdivisions
can prove yes on a found certificate and can fail to find any certificate
without proving no. Our argument supplies no uniform bound on the needed
subdivisions, no terminating negative-case test, and no resolution of the
even-genus surface case. In particular the vertex-map certificate alone
does not establish the complexity claimed in the broad catalogue question.

\section{Source check and adversarial audit}
The Matousek--Sedgwick--Tancer--Wagner primary abstract arXiv:1402.0815
reports decidability specifically for $\mathbb R^3$ or the 3-sphere. Its full
proof is not rederived here and the assertion is not transferred to an
arbitrary input 3-manifold. The audit checked the boundary-manifold case,
subdivision at the cone apex, and the separation between a polynomial
verifier on explicit refinements and a polynomial certificate bound for the
original input. All proved statements are auxiliary or rederivations;
no new complexity classification, computation or theorem-prover check is
asserted.

\section{Process and contributions}
Mathematics and model-assisted adversarial self-audit: OpenAI gpt-6-astra, ultra reasoning, as dispatched for this research batch; exact serving revision not exposed. Codex (GPT-6; exact serving revision not exposed) materially assisted with orchestration and format review. No human mathematical authorship, mathematical edits, or independent human proof checking is asserted.
The coordinating Codex (GPT-6; exact serving revision not exposed) also provided a separate model-assisted written-proof review of this manuscript. This is not independent human verification.
Contribution prompt A was applied to the full definition and supplied source row, then prompt B was applied adversarially to the derived results. The item-specific audit above records the scope checks and rejected extrapolations. Local assembly and structural checks do not establish mathematical correctness. The source catalogue status is preserved in TOP\_PROBLEM and is separate from this attempt status.
\paragraph{Integration review.} Codex (GPT-6; exact serving revision not exposed) reviewed the written arguments and expanded layout and source-reference wrappers. Repository metadata and the production web renderer were checked using Python 3.13 and JavaScript with MathJax 3.2.2. This model-assisted review is not independent human proof verification; no PDF compilation or formal proof checking is asserted.

\section{Completion Estimate}
\noindent\textbf{COMPLETION: 0\%}
The proved material is a rederivation of known or auxiliary facts rather than a substantiated new advance on the original target. The broad embedding and thickenability complexity questions, especially the even-genus and ambient-input cases, are not classified; local links and explicit refinement certificates do not close the gap.
\section{References}
The complete original Sources subsection above is retained, including catalogue attribution and licence notices. The additional list records primary sources actually consulted for this attempt on 11 October 2026; uncited source links reproduced above were not all rechecked.
\begin{itemize}
\item de Mesmay, Embeddability of 2-complexes, OWR 45/2015, pp. 2666--2668; original subquestions and cone obstruction read. \url{https://doi.org/10.4171/owr/2015/45}.
\item Matousek, Sedgwick, Tancer and Wagner, Embeddability in the 3-sphere is decidable; abstract inspected, general proof not independently checked. \url{https://arxiv.org/abs/1402.0815}.
\end{itemize}
\section{Final status}
\textbf{Unresolved in this attempt.} The broad embedding and thickenability complexity questions, especially the even-genus and ambient-input cases, are not classified; local links and explicit refinement certificates do not close the gap.
\end{document}
